%% cover_letter.tex — C13: CUEP Conservative Lower Bound
%% IEEE Transactions on Power Systems
\documentclass[11pt,a4paper]{letter}
\usepackage[T1]{fontenc}
\usepackage[utf8]{inputenc}
\usepackage[margin=2.5cm]{geometry}
\usepackage{microtype,amsmath}

\signature{Anoop~V.~Eluvathingal\\
           Ph.D.\ Candidate, Dept.\ of Electrical Engineering\\
           Indian Institute of Technology Madras\\
           Chennai 600036, India\\[2pt]
           \texttt{arun.eluvathingal@iitm.ac.in}}

\address{Department of Electrical Engineering\\
         Indian Institute of Technology Madras\\
         Chennai 600036, India}

\begin{document}
\begin{letter}{Editor-in-Chief\\
               \textit{IEEE Transactions on Power Systems}}

\opening{Dear Editor,}

We submit for your consideration the manuscript:

\begin{center}
  \textbf{Conservativeness of the CUEP Energy Bound in IBR-Penetrated
  Networks: Proof, Monotonicity, and Application to the SAMBPS
  Protection Framework}
\end{center}

\textbf{The problem.}
The CUEP energy bound is the foundation of direct transient stability
assessment.  Its conservativeness property --- that $V_{\rm cuep} \leq
V_{\rm cr}$, so the bound never declares an unstable clearing as stable
--- was proved for pure synchronous machine networks.  With inverter-based
resources (IBRs) penetrating distribution and sub-transmission networks
at 20--60\%, the IBR current limiter introduces a bounded non-gradient
force term into the swing equations, and the original proof does not
directly apply.  Whether the CUEP bound remains conservative under IBR
injection is an open question with direct implications for energy-function-based
protection relaying.

\textbf{Contributions.}
\begin{enumerate}
  \item \textbf{Conservativeness theorem (Theorem~1)}: We extend the
        Chiang--Hirsch--Wu proof to IBR-modified swing dynamics, showing
        $V_{\rm cuep}(k_{\rm ibr}) \leq V_{\rm cr}(k_{\rm ibr})$ for all $k_{\rm ibr} \leq 0.55$
        via a modified Lyapunov energy function that incorporates the
        IBR active power injection.

  \item \textbf{Monotonicity theorem (Theorem~2)}: We prove that the
        conservatism gap $\Delta V = V_{\rm cr} - V_{\rm cuep}$ is
        monotonically non-decreasing in $k_{\rm ibr}$ --- higher IBR penetration
        makes the bound \emph{safer}, not less safe.

  \item \textbf{SAMBPS C13 Proposition}: Zero false trips of the
        CUEP-based trip rule for $k_{\rm ibr} \leq 0.55$, with quantified
        CCT margin of 12--28\,ms (reducible to 6\,ms with an
        IBR-corrected CUEP).

  \item \textbf{IEEE 39-bus validation}: 1500 scenarios (5 penetration
        levels $\times$ 15 fault locations $\times$ 4 fault types);
        zero violations of Theorem~1 for $k_{\rm ibr} \leq 0.55$.
\end{enumerate}

\textbf{Fit to the journal.}
\textit{IEEE Transactions on Power Systems} is the primary venue for
rigorous stability theory and transient stability assessment.
This paper makes a new and complete theoretical contribution to direct
methods --- proving that a central property of the CUEP method survives
the transition to IBR-dominated grids --- and backs it with a large
simulation study.

\textbf{Suggested reviewers.}
\begin{itemize}
  \item Prof.\ Hsiao-Dong Chiang (Cornell University)
        --- CUEP method, direct stability
  \item Prof.\ Vijay Vittal (Arizona State University)
        --- transient stability with IBRs
  \item Prof.\ Thierry Van Cutsem (University of Li\`{e}ge)
        --- Lyapunov stability, power system dynamics
  \item Dr.\ Mani Venkatasubramanian (Washington State University)
        --- energy functions, stability boundaries
\end{itemize}

This manuscript is not under review elsewhere.
All authors have approved the submission.

\closing{Sincerely,}
\end{letter}
\end{document}
